% ===========================================================================
\section{Related Work}
\label{sec:related_work}
% ===========================================================================

\textbf{Blockchain provenance for AM.}
Traceability frameworks for AM record lifecycle events on permissioned
ledgers, typically Hyperledger Fabric~\cite{androulaki2018hyperledger}:
Alkhader et al.~\cite{alkhader2020blockchain} propose smart contracts for
AM part traceability, and Kurpjuweit et al.~\cite{kurpjuweit2021blockchain}
survey blockchain adoption across AM supply chains.  Our prior
work~\cite{anonymous2026unified} co-deploys provenance, reputation, and
integration contracts in one Fabric chaincode so lifecycle evidence and
actor trust scores update atomically, anchoring $\sim$337\,B of metadata
per event with an off-chain artifact hash.  All of these systems, ours
included, accept submitted event data at face value: the chaincode
validates \emph{who} submitted and \emph{whether the lifecycle transition
is legal}, not whether the reported physical process occurred.  This paper
addresses that residual gap.

\textbf{The oracle problem.}
The gap manifests the blockchain oracle
problem~\cite{caldarelli2020oracle}: smart contracts cannot verify
real-world facts, so external data carriers reintroduce the trusted third
party the ledger was meant to remove.  Decentralized oracle networks such
as Chainlink~\cite{breidenbach2021chainlink} mitigate this for
\emph{redundantly observable} data (e.g., market prices) by aggregating
independent reporters; manufacturing telemetry is not redundantly
observable---only the machine that printed the part witnessed its thermal
history---so the problem shifts to \emph{authenticating the single
witness}.  Town Crier~\cite{zhang2016towncrier} does this for web sources
using SGX~\cite{costan2016sgx} enclaves; our T3 tier
(Section~\ref{sec:tiers}) is the AM-floor analogue, while T1/T2 are
deployable intermediates needing no trusted hardware on legacy printers.

\textbf{AM security.}
Yampolskiy et al.~\cite{yampolskiy2018security} taxonomize AM attacks,
including sabotage via manipulated process parameters; the dr0wned
attack~\cite{belikovetsky2017dr0wned} demonstrated design-file sabotage
producing a drone propeller that failed in flight.  Such attacks motivate
binding the \emph{executed} design hash and \emph{observed} process
envelope into the provenance record.

\textbf{Reputation systems.}
Beta-distribution models~\cite{josang2002beta, josang2007survey} provide
constant-time Bayesian trust updates; EigenTrust~\cite{kamvar2003eigentrust}
popularized transitive peer trust.  On-chain reputation systems typically
depend on \emph{subjective} ratings; our bridge instead derives the rating
from measured physical adherence, grounding the signal in sensor evidence
whose integrity the attestation layer protects.  To our knowledge, no
prior AM-blockchain work implements and evaluates a device-signed,
canonically hashed telemetry attestation path on real printer hardware
with an explicit statement of its trust boundary.
